% DRAFT for the author's review. Not submitted. See paper/README.md.
\documentclass[11pt]{amsart}
\usepackage{amssymb,amsmath,amsthm}
\usepackage[margin=1.1in]{geometry}
\usepackage{booktabs}
\usepackage{hyperref}

\newtheorem{theorem}{Theorem}
\newtheorem{proposition}[theorem]{Proposition}
\newtheorem{lemma}[theorem]{Lemma}
\newtheorem{corollary}[theorem]{Corollary}
\theoremstyle{definition}
\newtheorem{definition}[theorem]{Definition}
\newtheorem{question}[theorem]{Question}
\theoremstyle{remark}
\newtheorem{remark}[theorem]{Remark}

\newcommand{\F}{\mathbb{F}}
\newcommand{\K}{\mathbb{K}}
\newcommand{\abs}[1]{\left|#1\right|}
\newcommand{\frp}[1]{\langle #1\rangle}

\title[Sharpness of the Lai--Sprang bound]{The Lai--Sprang lower bound for the
$P(t)$-adic Littlewood conjecture is attained in characteristics
$5,13,29,37,53$ and $61$}
\author{Greg Baker}
\date{Draft, 18 October 2026}

\begin{document}

\begin{abstract}
Lai and Sprang disproved the $P(t)$-adic Littlewood conjecture in every odd
characteristic $p$ with an explicit Laurent series $\Lambda$, proving
$\inf_{Q\neq0}\abs{Q}\,\abs{Q}_{P}\,\abs{\frp{Q\Lambda(P)}}\ge 2^{-2N\deg P}$,
where $N=2^{v_2(p-1)}$. We show that this constant is optimal when
$p\in\{5,13,29,37,53,61\}$: in each case an explicit multiplier attains
$2^{-2N}$ for $P=t$, and composition gives equality for every irreducible $P$
over every field of characteristic $p$. The witnesses come from a reformulation
of the extremal problem as the occurrence of a partial quotient of degree
$N+1$ in a shifted continued fraction, combined with a dyadic reverse lift.
All six cases are verified in the Lean proof assistant. In
particular, the sharper bound $2^{-N}$ suggested by small computations is
false. A random model predicts witnesses of degree about $p^{N/2}$; whether the
constant is attained for every odd $p$ remains open.
\end{abstract}

\maketitle

\section{Introduction}

Let $\K$ be a field and $\K((t^{-1}))$ the field of Laurent series in $t^{-1}$,
with $\abs{f}=2^{\deg f}$. For an irreducible $P\in\K[t]$ write
$\abs{Q}_P=\abs{P}^{-v_P(Q)}$, and let $\frp{f}$ denote the part of $f$ of
negative degree. The $P(t)$-adic Littlewood conjecture of de Mathan and
Teuli\'e asserts that
$\inf_{Q\in\K[t]\setminus\{0\}}\abs{Q}\,\abs{Q}_P\,\abs{\frp{Q\Theta}}=0$
for every $\Theta\in\K((t^{-1}))$. It fails over infinite fields; over finite
fields counterexamples were found in characteristic $3$ by Adiceam, Nesharim
and Lunnon~\cite{ANL}, in characteristics $5$, $7$, $11$ by Garrett and
Robertson~\cite{GR}, and in every odd characteristic by Lai and
Sprang~\cite{LS}. Characteristic two remains open.

Fix an odd prime $p$, let $r=v_2(p-1)$, $N=2^r$ and $H=N/2$. Lai and Sprang
consider
\[
\Lambda(t)=\sum_{k\ge0}\ \sum_{\zeta^{H}=-1}\frac{t^{2^k}}{t^{2^{k+1}}-\zeta}
\in\F_p((t^{-1})),
\]
the inner sum running over the $H$ primitive $N$-th roots of unity in $\F_p$,
and prove \cite[Theorem~1.2]{LS}
\begin{equation}\label{eq:LS}
\inf_{Q\in\K[t]\setminus\{0\}}\abs{Q}\,\abs{Q}_P\,\abs{\frp{Q\Lambda(P)}}
\ \ge\ 2^{-2N\deg P}
\end{equation}
for every field $\K$ of characteristic $p$ and every irreducible $P\in\K[t]$.
They do not discuss optimality of the constant.

\begin{theorem}\label{thm:main}
Let $p\in\{5,13,29,37,53,61\}$, so that $N=4$. There is $Q\in\F_p[t]$ with
$\abs{Q}\,\abs{Q}_t\,\abs{\frp{Q\Lambda}}=2^{-8}$. Consequently equality holds
in \eqref{eq:LS} for every field $\K$ of characteristic $p$ and every
irreducible $P\in\K[t]$.
\end{theorem}

For $p=5$ one may take $Q=t^{1244}S(t^4)$ with
\[
S=(1+t+t^2+t^3)(2+t+t^2+t^3+3t^4+3t^5+t^6+t^7+t^8),
\]
so that $S(t^4)$ has degree $44$ and $Q\Lambda$ has its first nonzero
fractional coefficient at $t^{-52}$. Every $N=4$ prime up to $61$ is covered.
Small computations had suggested the sharper bound $2^{-N}$: it holds for all
multipliers of degree below $44$ when $p=5$ (Theorem~\ref{thm:minimal}), but
not in general.

The proof of Theorem~\ref{thm:main} is a finite computation together with a
reverse lift (Lemma~\ref{lem:lift}) that turns a short ``terminal'' window into
an extremal multiplier. The witnesses were found by recasting the existence of
a terminal window as a continued-fraction question (Section~\ref{sec:cf}).

\section{Coefficients, rows and the reverse lift}

Write $\Lambda=\sum_{n\ge1}a_nt^{-n}$. Expanding
$t^{2^k}/(t^{2^{k+1}}-\zeta)=\sum_{u\ge0}\zeta^u t^{-2^k(2u+1)}$ gives
$a_{2^k(2u+1)}=\sum_\zeta\zeta^u$. The roots form a coset of the $H$-th roots of
unity, so $\sum_\zeta\zeta^u=H(-1)^{u/H}$ if $H\mid u$ and $0$ otherwise. Hence
\begin{equation}\label{eq:coeff}
a_{2n}=a_n,\qquad a_n=H(-1)^h\ \text{if the odd part of $n$ is $1+Nh$},
\qquad a_n=0\ \text{otherwise}.
\end{equation}
For $T=\sum_iT_it^i\in\F_p[t]$ and $v\ge1$ set $B_T(v)=\sum_iT_ia_{v+i}$, the
coefficient of $t^{-k}$ in $t^MT\Lambda$ when $v=M+k$. If $Q=t^mR$ with
$R(0)\neq0$ and $\deg R=d$, and $j\ge1$ is least with $B_R(m+j)\neq0$, then
\begin{equation}\label{eq:product}
\abs{Q}\,\abs{Q}_t\,\abs{\frp{Q\Lambda}}=2^{m+d}\,2^{-m}\,2^{-j}=2^{d-j}.
\end{equation}
We call $j-d$ the \emph{defect}. By \eqref{eq:LS} every defect is at most
$2N$.

\begin{lemma}[reverse lift]\label{lem:lift}
Let $F=1+t+\dots+t^{N-1}$ and suppose $F\mid S$. Then for all $v\ge1$,
$B_{S(t^N)}(v)=B_S(v/N)$ if $N\mid v$, and $B_{S(t^N)}(v)=0$ otherwise.
\end{lemma}

\begin{proof}
Put $T_i=S(t^{2^i})$ for $0\le i\le r$. Let $1\le i\le r$. By \eqref{eq:coeff},
$B_{T_i}(2w)=\sum_lT_{i-1,l}\,a_{2(l+w)}=B_{T_{i-1}}(w)$. For odd arguments,
\[
B_{T_i}(2w+1)=\sum_lT_{i-1,l}\,a_{2(l+w)+1}
=\sum_\zeta\zeta^w\,T_{i-1}(\zeta)=\sum_\zeta\zeta^w\,S\bigl(\zeta^{2^{i-1}}\bigr).
\]
Each $\zeta^{2^{i-1}}$ is an $N$-th root of unity of order $N/2^{i-1}\ge2$, so
it is a root of $F$, and every summand vanishes. Induction on $i$ gives the
claim for $T_r=S(t^N)$.
\end{proof}

\begin{definition}
A \emph{terminal window} is a pair $(S,M)$ with $F\mid S$, $S\neq0$,
$D=\deg S$, $M\ge0$, and $B_S(M+k)=0$ for $1\le k\le D+1$.
\end{definition}

\begin{corollary}\label{cor:lift}
Let $(S,M)$ be a terminal window with $S(0)\neq0$ and $B_S(M+D+2)\neq0$. Then
$Q=t^{NM}S(t^N)$ has $\abs{Q}\,\abs{Q}_t\,\abs{\frp{Q\Lambda}}=2^{-2N}$.
\end{corollary}

\begin{proof}
$R=S(t^N)$ has $R(0)=S(0)\neq0$ and degree $ND$. By Lemma~\ref{lem:lift} its
first nonzero row at shift $NM$ is row $N(D+2)$. Apply \eqref{eq:product}.
\end{proof}

\begin{remark}\label{rem:cap}
The hypothesis $B_S(M+D+2)\neq0$ is automatic: Lai--Sprang's
Proposition~2.3~\cite{LS} applied to $(t-1)S$, a multiple of $t^N-1$, shows
that no $F$-multiple has $D+2$ zero rows. For the same reason every terminal
window has $S(0)\neq0$. We verify both
conditions directly, so the proof of Theorem~\ref{thm:main} does not depend on
this remark.
\end{remark}

\section{Proof of Theorem~\ref{thm:main}}\label{sec:proof}

Table~\ref{tab:witnesses} lists, for each prime, an odd shift $M$ and a
polynomial $S=FV$ with $\deg V=e$. In each case $S(0)\neq0$, the rows
$B_S(M+1),\dots,B_S(M+D+1)$ vanish and $B_S(M+D+2)\neq0$. These are finite
statements about at most $2D+M+2$ coefficients given by \eqref{eq:coeff}.
The polynomial for $p=5$ is displayed above; the others are in the ancillary
data~\cite{repo}. Corollary~\ref{cor:lift} gives a multiplier of degree $4D$
with defect $8=2N$, proving the first assertion.

\begin{table}[ht]
\centering
\begin{tabular}{rrrrrrr}
\toprule
$p$ & $M$ & $e=\deg V$ & $D=\deg S$ & $\deg R=4D$ & shift $4M$ & $j=4(D+2)$\\
\midrule
5 & 311 & 8 & 11 & 44 & 1244 & 52\\
13 & 1465 & 76 & 79 & 316 & 5860 & 324\\
29 & 2501 & 326 & 329 & 1316 & 10004 & 1324\\
37 & 19939 & 1134 & 1137 & 4548 & 79756 & 4556\\
53 & 40193 & 1128 & 1131 & 4524 & 160772 & 4532\\
61 & 22815 & 2083 & 2086 & 8344 & 91260 & 8352\\
\bottomrule
\end{tabular}
\caption{Terminal windows and their lifts, $N=4$.}\label{tab:witnesses}
\end{table}

For the second assertion let $Q=t^mR$ be as above, $d=\deg R$, and let
$\frp{Q\Lambda}=f=\sum_{k\ge j}f_kt^{-k}$, $f_j\neq0$. Substituting $t\mapsto P$
keeps the polynomial part of $Q\Lambda$ polynomial, while
$f(P)=\sum_{k\ge j}f_kP^{-k}$ has absolute value $2^{-j\deg P}<1$; hence
$\frp{Q(P)\Lambda(P)}=f(P)$. Also $\abs{Q(P)}=2^{(m+d)\deg P}$ and, as
$R(P)\equiv R(0)\neq0\pmod P$, $v_P(Q(P))=m$. The product for $Q(P)$ is
therefore $2^{(d-j)\deg P}=2^{-2N\deg P}$, which equals the bound
\eqref{eq:LS}. (Compare Robertson's transfer principle~\cite{Rob}.)
\qed

\begin{remark}[formal verification]
For $p=5$ the exact certificate of $Q$ and the identity
$\abs{Q}\abs{Q}_t\abs{\frp{Q\Lambda}}=2^{-8}<2^{-4}$ are checked by the Lean
kernel for the actual polynomial and the infinite coefficient stream
\eqref{eq:coeff}. For the other primes the Lean proofs check the terminal
certificate and derive the lifted certificate from a formal version of
Lemma~\ref{lem:lift}: for any stream with $a_{2n}=a_n$, vanishing odd rows
double the shift and the first index. The odd-row condition reduces to finitely
many residues by periodicity of $a_{2u+1}$ in $u$. For $p\ge29$ the terminal
rows are evaluated by a natural-number checker with a soundness proof. All
proofs use only the standard axioms; see \cite{repo}.
\end{remark}

\section{A continued-fraction formulation}\label{sec:cf}

Let $b_v=a_v+\dots+a_{v+N-1}$, the coefficients of $F\Lambda$. For $S=FV$ the
rows of $S$ at shift $M$ are $\sum_jV_jb_{M+k+j}$. Thus $(FV,M)$ is a terminal
window with $\deg V=e$ exactly when
$\abs{\frp{V\,\Theta_M}}\le2^{-(N+e+1)}$, where
\[
\Theta_M=\frp{t^MF\Lambda}=\sum_{k\ge1}b_{M+k}t^{-k}.
\]
By the best-approximation property of continued fractions this happens if and
only if $\Theta_M$ has a convergent denominator of degree at most $e$ whose next
partial quotient has degree at least $N+1$. Remark~\ref{rem:cap} shows that
no partial quotient of any $\Theta_M$ has degree above $N+1$.

\begin{proposition}\label{prop:cf}
Let $\Delta_p$ be the largest defect of a nonzero multiplier at any shift.
Then $\Delta_p\in\{N,2N\}$, and $\Delta_p=2N$ if and only if some
$\Theta_M$ has a partial quotient of degree $N+1$.
\end{proposition}

\begin{proof}[Sketch]
$\Delta_p\ge N$ is attained in degree zero by $R=1$, $m=5N+1$. If some
partial quotient of $\Theta_M$ has degree $N+1$, Corollary~\ref{cor:lift}
gives defect $2N$. Conversely, a multiplier with defect above $N$ descends, by
the factor-aware parity descent adapted from~\cite[\S3]{LS}, to a terminal
window; this ordinary argument is in~\cite{repo}, where its descent steps are
also formalised.
\end{proof}

A single Euclidean expansion of $\Theta_M$ truncated after $2E+N+1$ terms
decides every quotient degree $e\le E$ at that shift. All shifts reduce to a
finite list: put $T=2N+2E$, let $L$ be the least power of two exceeding $T$
and write $M=qNL+b$. By \eqref{eq:coeff}, every coefficient $a_{M+i}$,
$1\le i\le T$, equals $a_{b+i}$ unless $L\mid b+i$, which happens for at
most one $i$; there the value lies in $\{0,\pm H\}$. Running the expansion
for the $NL$ residues $b$, with all three values at the exceptional index,
therefore covers every shift.

\begin{theorem}\label{thm:minimal}
For $p=5,13,29,37,53,61$ every nonzero $R$ of degree $d$ has a nonzero row
among the first $d+N$ rows, at every shift, whenever
$d<44,316,1316,4548,4524,8344$ respectively. Each threshold is attained by Table~\ref{tab:witnesses}.
\end{theorem}

This is computer-assisted: the covers exclude terminal windows of smaller
quotient degree, and the descent transfers the exclusion to all multipliers.
Two independent implementations (Berlekamp--Massey linear complexity in Python
and Euclidean division in C) give identical results where both were run, and
for $p=5$ a direct nullspace computation agrees.

The integer matrix of the $p=5$ terminal window has the greatest common divisor
of its maximal minors equal to $5$. So the witness is specific to
characteristic $5$, not an integer rank drop.

\section{Heuristics and questions}

Model each realizable window and degree as an independent event of
probability $p^{-N}$ for a zero block of side $N$ in the number wall of $b$
modulo $p$. Because the sequence is $2$-automatic, distinct windows up to
degree $E$ number on the order of $E^2$, so the model predicts least terminal
degrees of order $p^{N/2}$. For $N=4$ the observed least degrees
$8,76,326,1134,1128,2083$ are $0.32,0.45,0.39,0.83,0.40,0.56$ times $p^2$. For $N=8$ the model predicts
degrees around $3\cdot10^6$ already at $p=41$. Complete covers to quotient
degree $250$ found nothing for $p=41,73,89,137$ ($N=8$) and $p=17,113$
($N=16$), which is consistent with the model and is no evidence for the
sharper constant.

\begin{question}
Is $\inf_Q\abs{Q}\abs{Q}_t\abs{\frp{Q\Lambda}}=2^{-2N}$ for every odd prime
$p$? For $r=1$ this holds since $R=1+t^2$, $m=2$ has rows $0,0,0,0,0,2$.
\end{question}

\begin{remark}
Garrett and Robertson~\cite[Conjecture~1.8]{GR} predict, for their level-$n$
paperfolding series, the infimum $q^{-2^n}$ with $2^n=N$. As printed this
contradicts~\cite[Theorem~1.1]{ANL}, which gives deficiency $4=2^{n+1}$ in
characteristic $3$, so the intended exponent is presumably $2^{n+1}$. For
their level-two series over $\F_5$ the multiplier $t^{61}R$ with $\deg R=64$
has first index $72$, deficiency $8=2^{n+1}$. Scans to shift $60{,}000$ found
no larger value. This agrees with the reading $2^{n+1}$ and with
Theorem~\ref{thm:main} for the Lai--Sprang series.
\end{remark}

\subsection*{Data, code and use of AI}
All witnesses, both cover implementations, independent root-sum readback and
the Lean proofs are in the public repository~\cite{repo}. The computations,
Lean proofs and a first draft of this note were produced with substantial
assistance from AI coding agents (Anthropic Claude Code; earlier repository
work used OpenAI Codex), under the author's direction. Every claim is checked
by an exact finite computation, a kernel-checked proof, or an ordinary proof
in the repository's dated reports.

\begin{thebibliography}{9}
\bibitem{ANL} F.~Adiceam, E.~Nesharim and F.~Lunnon, On the $t$-adic
Littlewood conjecture, \emph{Duke Math.\ J.} \textbf{170} (2021), 2371--2419;
arXiv:1806.04478.
\bibitem{GR} S.~Garrett and S.~Robertson, Counterexamples to the $p(t)$-adic
Littlewood conjecture over small finite fields, arXiv:2405.14454v2 (2025).
\bibitem{LS} L.~Lai and J.~Sprang, On the $P(t)$-adic Littlewood conjecture in
odd characteristics, arXiv:2606.00633v1 (2026).
\bibitem{Rob} S.~Robertson, Combinatorics on number walls and the
$P(t)$-adic Littlewood conjecture, arXiv:2307.00955v3 (2025).
\bibitem{repo} G.~Baker, Function-field $P(t)$-adic Littlewood research,
\url{https://github.com/solresol/function-field-littlewood}, files
\texttt{results/2026-10-17-lai-sprang-sharpness.md},
\texttt{FunctionFieldLittlewood/Sharpness.lean},
\texttt{FunctionFieldLittlewood/ReverseLift.lean}.
\end{thebibliography}

\end{document}
